\section{A sharp detecting scale when a double endpoint controls the angle}
\label{sec:analytic-dominant-double}

The following conditional case closes a matching upper-bound gate for
the actual heat family. The extra multiplicity hypothesis is displayed;
it is not part of A. We first prove the infinite-alphabet estimate that
makes the negative minor uniform near a coalescing pair.

\begin{lemma}[A coalescing dominant pair in rectangular Schur minors]
\label{lem:analytic-rectangular-pair}
Let an $\ell^1$ complex alphabet, closed under conjugation, depend on
a positive parameter tending to zero. Suppose its entries consist of
\[
 b_1>\cdots>b_p>b,\qquad
 \beta_+=be^{i\theta},\quad\beta_-=be^{-i\theta},\qquad
                         (\eta_j)_{j\ge1},
\]
where $\theta\downarrow0$. Assume $b_i\to b_i^0$,
$b\to b^0>0$, with strictly separated limits
$b_1^0>\cdots>b_p^0>b^0$ when $p>0$ (vacuous for $p=0$), and
$\sup_j|\eta_j|\le c_0<b^0$ uniformly. Assume also
$\sum_j|\eta_j|$ is uniformly bounded, and that the tail products
\[
             C(z)=\prod_j(1-\eta_jz)^{-1}
\]
converge locally uniformly on $|z|<1/c_0$ to the corresponding product of an endpoint
alphabet of positive real entries less than $b^0$. Write $k=p+1$,
$P=\prod_{i=1}^p b_i$, and $h_r$ for the complete symmetric functions
of the full alphabet. Then there are constants $C_1>0$, $0<q<1$,
independent of $n$ and of the sufficiently small parameter, such that
\begin{equation}\label{eq:analytic-rectangular-asymptotic}
 \frac{s_{(n^k)}}{(Pb)^n}
  =\frac{|D|\sin((n+1)\theta+\arg D)}{\sin\theta}
                         +O(C_1\theta^{-k}q^n),\qquad n\ge k,
\end{equation}
where
\[
 D=\left\{\prod_{i=1}^p
              \frac{C(1/b_i)}{1-\beta_-/b_i}\right\}C(1/\beta_+)
                             \longrightarrow D_0>0.
\]
In particular if $n\theta\to u\in(\pi,2\pi)$, the displayed
rectangular Schur value is strictly negative eventually.
The same value is the elementary Toeplitz minor
\begin{equation}\label{eq:analytic-rectangular-elementary-minor}
                  s_{(n^k)}=det[e_{k-i+j}]_{i,j=1}^n.
\end{equation}
It has order $n$ and fits in the index window $0,\ldots,k+n-1$.
\end{lemma}

\begin{proof}
Choose $c$ strictly between $c_0$ and $b^0$, and smaller than
all the limiting $b_i$. For sufficiently small parameter, the generating
function
\[
 H(z)=\sum_{r\ge0}h_rz^r
       =\prod_{i=1}^p(1-b_i z)^{-1}
           (1-\beta_+z)^{-1}(1-\beta_-z)^{-1} C(z)
\]
has exactly these $p+2$ simple poles inside $|z|=1/c$.
Its modulus on that circle is uniformly bounded. Remove the simple
partial fractions there:
\[
 H(z)=\sum_{\alpha\in\mathcal B}
          \frac{A_\alpha}{1-\alpha z}+G(z),\qquad
 \mathcal B=\{b_1,\ldots,b_p,\beta_+,\beta_-\},
\]
where
\[
 A_\alpha=\prod_{\gamma\ne\alpha}(1-\gamma/\alpha)^{-1}.
\]
The product includes the whole tail. The tail product converges
absolutely and is bounded above and away from zero at the required
points, by its uniform $\ell^1$ bound and strict modulus separation.
The only vanishing denominator is
$1-\beta_-/\beta_+$ or its conjugate. Hence
$|A_\alpha|\le C\theta^{-1}$, and Cauchy's formula gives
\begin{equation}\label{eq:analytic-h-poles-remainder}
 h_r=\sum_{\alpha\in\mathcal B}A_\alpha\alpha^r+g_r,
                         \qquad |g_r|\le C\theta^{-1}c^r.
\end{equation}
The remainder $G$ is analytic through all removed poles. Its circle
bound follows from those of $H$ and the partial fractions, so no
uniform separation of the two coalescing poles was assumed.

Use the ordinary Jacobi--Trudi determinant
$s_{(n^k)}=\det[h_{n-i+j}]_{i,j=0}^{k-1}$.
The pole part of its matrix is a sum of rank-one matrices
\[
        A_\alpha\alpha^n
              (\alpha^{-i})_{i=0}^{k-1}
              (\alpha^{j})_{j=0}^{k-1}.
\]
Cauchy--Binet for that finite sum gives one term for each
$k$-element subset $J\subset\mathcal B$. The two Vandermonde
determinants cancel all internal factors in
$\prod_{\alpha\in J}A_\alpha$, leaving exactly
\begin{equation}\label{eq:analytic-rectangular-subset-coefficient}
 \left(\prod_{\alpha\in J}\alpha\right)^n
       \prod_{\substack{\alpha\in J\\\gamma\notin J}}
                                      (1-\gamma/\alpha)^{-1}.
\end{equation}
Here the complement includes the infinite tail, only through its
absolutely convergent products. To see the cancellation directly,
each pair $\alpha,\alpha'$ within $J$ contributes
$-(\alpha\alpha')/(\alpha-\alpha')^2$ from the two partial
fractions, the inverse of its contribution from the two Vandermondes.

The two subsets containing every $b_i$ and exactly one of
$\beta_\pm$ have modulus weight $(Pb)^n$. Their coefficients are
$D/(1-e^{-2i\theta})$ and its conjugate. Since
\[
             (1-e^{-2i\theta})^{-1}
                         =e^{i\theta}/(2i\sin\theta),
\]
their sum is precisely the first term of
\eqref{eq:analytic-rectangular-asymptotic}. The endpoint tail is
positive real, all $b_i^0>b^0$, and the displayed products therefore
give $D\to D_0>0$.

Every other $k$-element subset omits at least one $b_i$ and
contains both members of the pair. Its modulus weight divided by
$Pb$ is $b/b_i$, bounded by a constant less than one. Its coefficient
is uniformly bounded: the coalescing pair is internal to $J$, so its
potential singular factors already cancelled in
\eqref{eq:analytic-rectangular-subset-coefficient}.

It remains to control determinants that use remainder columns from
\eqref{eq:analytic-h-poles-remainder}. This cannot be done by a
naive bound using the largest pole to the $k$th power. Instead use
multilinearity. A term with $m\ge1$ remainder columns has $k-m$
remaining rank-one pole columns, whose pole labels must be distinct
or the determinant is zero. Its modulus is at most a fixed constant
 times
\[
 \theta^{-k}\left(c^m\prod_{\alpha\in J}|\alpha|\right)^n,
                              \qquad |J|=k-m.
\]
Fixed shifts $j-i$ affect only the constant. The largest product of
$k-m\le p$ pole moduli uses that many of the $b_i$. Because each
remaining $b_i$ is larger than $b$,
\[
          \frac{c^m\prod_{\alpha\in J}|\alpha|}{Pb}
                                      \le(c/b)^m\le q_1<1.
\]
There are finitely many such multilinear terms, their number depending
only on $k$. Combining this bound with the nonleading subsets proves
\eqref{eq:analytic-rectangular-asymptotic}, after increasing $q<1$.
If $n\theta\to u\in(\pi,2\pi)$, its first term is asymptotic to
$D_0\sin u/\theta<0$, and its exponential remainder is smaller than
any power of $\theta$. Finally dual Jacobi--Trudi transposes
$(n^k)$ to $(k^n)$ and gives
\eqref{eq:analytic-rectangular-elementary-minor}. Explicitly use
row indices $0,\ldots,n-1$ and column indices $k,\ldots,k+n-1$
in the upper-triangular Toeplitz matrix $(e_{j-i})$.
\end{proof}

\begin{theorem}[NN-A027: exact detecting constants in the dominant-double case]
\label{thm:analytic-dominant-double}
Assume A. For a positive zero $x$ of $H_0$, let $m_x$ be its
multiplicity and $h_{m_x}$ the largest absolute root of the
probabilists' Hermite polynomial $\operatorname{He}_{m_x}$,
with $h_1=0$. Make the following additional hypothesis:
there is a positive \emph{double} zero $x_0$ such that
\begin{equation}\label{eq:analytic-dominant-double-hypothesis}
                    \frac{h_{m_x}}x\le\frac1{x_0}
                    \qquad\hbox{for every positive zero }x.
\end{equation}
This is a precise conditional case; neither this multiplicity pattern
nor existence of any multiple zero is inferred from A.

Let $m_*(-a)$ be the least order of a negative Toeplitz minor of
$F_{-a}$, and $R_*(-a)$ its least detecting index window as in
NN-A014 and NN-A020. Then
\begin{equation}\label{eq:analytic-sharp-order-window}
       \lim_{a\downarrow0}\sqrt a\,m_*(-a)
        =\lim_{a\downarrow0}\sqrt a\,R_*(-a)
                       =\frac{\pi x_0}{2\sqrt2}.
\end{equation}
For the normalized unitary characters of NN-A021, put
\[
 p_0=F_0'(0)=\sum_{x>0}x^{-2},\qquad
 \beta_0=\frac1{p_0x_0^2}\in(0,1),
\]
and let $d_*(-a)$ be the least rank at which
$U\mapsto\det E_{-a}(U-I)$ fails to be positive type. Then
\begin{equation}\label{eq:analytic-sharp-unitary-rank}
       \lim_{a\downarrow0}\sqrt a\,d_*(-a)
                       =\frac{\pi x_0(1-\beta_0)}{2\sqrt2}.
\end{equation}
There are explicit negative rectangular witnesses: if $p$ is the
number of positive endpoint zeros below $x_0$ and $k=p+1$, take
$n\theta(a)\to u\in(\pi,2\pi)$, where $\theta(a)$ is the positive
argument of the corresponding complex product parameter. The minor
$\det[e_{k-i+j}]_{i,j=1}^n$ is negative. For the translated unitary
alphabet this is the negative Fourier coefficient of the actual
one-dimensional representation $\det(U)^k$ on $U(n)$.
\end{theorem}

\begin{proof}
For $m\ge2$, $h_m\ge h_2=1$, by the principal Jacobi-matrix
interlacing of the Hermite recurrence. Therefore
\eqref{eq:analytic-dominant-double-hypothesis} forces every endpoint
zero below $x_0$ to be simple. There are finitely many of them,
$x_1<\cdots<x_p<x_0$. Their continuations at time $-a$ remain
simple and real for small $a$. The double zero at $x_0$ splits by
NN-A006 into
\[
             z_\pm(a)=x_0\pm i\sqrt{2a}+O(a).
\]
The two parameters $z_\pm(a)^{-2}$ are conjugate and have arguments
$\pm\theta(a)$ with
\begin{equation}\label{eq:analytic-pair-theta-sharp}
                         \frac{\theta(a)}{\sqrt a}
                                      \longrightarrow\frac{2\sqrt2}{x_0}.
\end{equation}
The $p$ real parameters from lower zeros have strictly larger moduli.
All remaining parameters have uniformly smaller modulus than
$x_0^{-2}$: choose a circle between $x_0$ and the next larger positive
endpoint zero, and use compact zero convergence. Zeros outside a
large fixed circle contribute parameters tending uniformly to zero.

Here and below the infinite products and their uniform bounds are
legitimate without simplicity of the other zeros. For $t$ in a compact
interval the theta integral gives
$\log\max_{|z|\le R}|H_t(z)|\le C R\log(R+3)$, with $C$ uniform,
and $H_t(0)$ is bounded away from zero near $t=0$. Jensen's formula
therefore gives a uniform zero count $N_t(R)\le C R\log(R+3)$.
Partial summation bounds uniformly the tails of
$\sum|z_j(t)|^{-2}$ and makes them tend to zero as the radius tends
to infinity. The genus-zero products of $F_t$, after removing the
fixed finite cluster, consequently converge locally uniformly to the
endpoint tail product. In particular the hypotheses of
Lemma~\ref{lem:analytic-rectangular-pair} hold for the alphabet
$z_j(-a)^{-2}$, with $k=p+1$.

Take $n=\lceil(\pi+\varepsilon)/\theta(a)\rceil$ for any fixed
$0<\varepsilon<\pi$. The lemma gives a negative minor of order $n$
in window $k+n-1$ for all sufficiently small $a$. Hence
\[
 m_*(-a)\le n,\qquad R_*(-a)\le k+n-1,
\]
and \eqref{eq:analytic-pair-theta-sharp} gives the corresponding
upper limits with $\pi+\varepsilon$ in place of $\pi$.

For the lower limits, NN-A022 and the extra hypothesis give
\[
 \frac{\alpha_H(-a)}{\sqrt a}\longrightarrow\frac{\sqrt2}{x_0}.
\]
The maximal parameter angle is exactly $2\alpha_H(-a)$ when these
angles are small. The strict quadratic sector proof in NN-A020
shows that every allowed minor of order at most $r$ is positive
whenever $(r+1)2\alpha_H(-a)<\pi$. Thus for any fixed
$\varepsilon>0$ one may take
$r=\lfloor(\pi-\varepsilon)/(2\alpha_H(-a))\rfloor-1$.
It follows that
\[
 \liminf_{a\downarrow0}\sqrt a\,m_*(-a)
                           \ge\frac{\pi x_0}{2\sqrt2}.
\]
Any minor in the window $0,\ldots,R$ has order at most $R+1$,
so $R_*(-a)\ge m_*(-a)-1$. Letting $\varepsilon\downarrow0$
in the upper bounds proves \eqref{eq:analytic-sharp-order-window}.

For the unitary statement put $p_a=F_{-a}'(0)>0$. The translated
alphabet in NN-A021 is
\[
 \widetilde\beta_j(a)
   =\frac{\beta_j(a)}{1-\beta_j(a)}
   =\frac1{p_a z_j(-a)^2-1},\qquad
                    \beta_j(a)=\frac1{p_a z_j(-a)^2}.
\]
At the endpoint all its entries are positive and are strictly decreasing
as functions of the positive zero $x$. Also $\sup_j|\beta_j(a)|<b<1$
for small $a$, by the compact zero separation already proved in
NN-A021. Its $\ell^1$ tail and local convergence therefore satisfy
the same lemma. At an endpoint zero $x$ of multiplicity $m_x$,
the largest local angle coefficient becomes
\begin{equation}\label{eq:analytic-translated-angle-coefficient}
 \frac{2\sqrt2\,h_{m_x}}{x(1-1/(p_0x^2))}.
\end{equation}
This follows by substituting the local Hermite expansion into
$(p_a z^2-1)^{-1}$; $p_a=p_0+O(a)$ is real.

The maximum of these coefficients is attained at $x_0$.
Indeed the nonzero coefficients occur only at $x\ge x_0$,
$h_{m_x}/x\le1/x_0$, and
$x\mapsto(1-1/(p_0x^2))^{-1}$ decreases there. The uniform
radial-tail estimate in NN-A022 also applies after this transformation:
for high zeros $|\beta_j(a)|\le C/|z_j(-a)|^2$, while
\[
 \left|\arg\frac{\beta}{1-\beta}\right|
                         \le\frac{|\arg\beta|}{1-|\beta|}
                         \qquad(|\beta|<1).
\]
One proves this last bound by integrating
$\operatorname{Re}(1-re^{i\theta})^{-1}\le(1-r)^{-1}$
along the argument, starting at a positive real number.
Consequently the maximal translated angle, divided by $\sqrt a$,
has the exact limit
$2\sqrt2/[x_0(1-\beta_0)]$.

The sector lower bound and the same rectangular upper bound now give
\eqref{eq:analytic-sharp-unitary-rank}, using the exact $PF_d$
versus $U(d)$ equivalence of NN-A021. Multiplication by the positive
constant $\Psi_a(0)=E_{-a}(-1)$ does not affect minor signs.
For the upper witness, the coefficient formula from that equivalence
is
\[
 c_{(k^n),n}=\det[a_{k-i+j}]_{i,j=1}^n
             =\Psi_a(0)^n s_{(n^k)}(\widetilde\beta(a))<0.
\]
The highest weight $(k,\ldots,k)$ is precisely $\det(U)^k$,
so this is the claimed concrete irreducible witness.
\end{proof}

\paragraph{Geometric and representation meaning.}
The additional hypothesis states that a local Artin algebra of length
two controls the maximal scaled angular motion of the endpoint divisor.
Its quadratic normalization produces the two conjugate dominant
parameters. The rectangular Schur computation cancels the internal
Vandermonde factors and isolates their sine oscillation with a uniformly
controlled infinite complementary divisor. Dual Jacobi--Trudi turns
this into a finite minor whose order and index window have the exact
asymptotic above. After the actual unitary translation, that same
rectangle is a negative coefficient of $\det^k$ at rank $n$.
The exterior-power/representation-ring dictionary is SGA~6, Expos\'e~V,
Example~2.2(b) and Lemma~2.2.1; the rank-two discriminant criterion is
SGA~1, Expos\'e~I, section~4, Proposition~10. These are sharp bounds
for the specified analytic and representation certificate families
under the displayed multiplicity hypothesis. They are not a proof
of that hypothesis, a pure motive, or a lower bound against all
algorithms.
